\documentclass{article}

% --- Track option: Workshop (anonymous submission) ---
% No trained-model results exist; this is a reference architecture
% and validation-harness paper. Use submission mode (no option).
\usepackage{neurips_2026}

% --- Required packages ---
\usepackage[utf8]{inputenc}
\usepackage[T1]{fontenc}
\usepackage{hyperref}
\usepackage{url}
\usepackage{booktabs}
\usepackage{amsfonts}
\usepackage{amsmath}
\usepackage{amssymb}
\usepackage{nicefrac}
\usepackage{microtype}
\usepackage{xcolor}
\usepackage{graphicx}
% \usepackage{subcaption}   % not used in this version
% \usepackage{algorithm}     % not used in this version
% \usepackage{algpseudocode} % not used in this version

% --- Title & authors (anonymous submission) ---
\title{Action-Incremental DQN: A Reference Architecture and Validation Harness for Unexpected Action-Space Expansion in Deep Reinforcement Learning}

\author{%
  Anonymous Author(s) \\
  Affiliation \\
  Address \\
  \texttt{email} \\
}

\begin{document}

\maketitle

\begin{abstract}
Standard deep reinforcement learning assumes a fixed action space that is known before training and never changes. In many real-world settings---video games that unlock new abilities at higher levels, robots fitted with new end-effectors mid-deployment, and agents that discover new affordances---unexpected actions become available mid-training, and retraining from scratch discards prior knowledge. We present \textbf{Action-Incremental DQN}, a reference architecture and validation harness for this understudied problem, built on three mechanisms: (i) a factorized Q-function $Q(s,a) = f_Q(\varphi(s), e_a)$ that decouples action identity from value estimation via a learnable action-embedding table; (ii) $k$-nearest-neighbour centroid interpolation in embedding space that initialises new actions from known action structure; and (iii) a progressive three-phase unfreezing schedule designed to mitigate catastrophic forgetting during fine-tuning of the new action. The validation harness spans four layers---shape correctness (12 tests), gradient flow (5 tests), RL-specific properties (10 tests), and numerical stability including bf16 safety (9 tests)---supplemented by 17 synthetic micro-benchmarks (53 automated tests total). A 14-condition ablation study characterises the design space along embedding dimension, initialisation strategy, freeze schedule, Q-function form, and auxiliary loss configuration, using synthetic proxy metrics on random states. The custom Gymnasium benchmark environment, trained-model evaluation on real tasks, and baseline comparisons are explicitly out of scope and identified as natural follow-up work. All code is included in the accompanying artifacts as a reproducible starting point for the community.
\end{abstract}


\section{Introduction}

Deep reinforcement learning has achieved remarkable successes across domains from game play~\cite{mnih2015human} to continuous control~\cite{lillicrap2015continuous} and policy optimisation~\cite{schulman2017proximal}. A tacit assumption shared by nearly all these methods is that the action set $\mathcal{A}$ is fixed at the start of training and remains unchanged throughout the agent's lifetime. The agent's policy $\pi(a \mid s)$ and value function $Q(s,a)$ are parameterised for exactly $|\mathcal{A}|$ outputs, and any addition or removal of actions requires architectural modification and, typically, retraining from scratch.

This assumption fails in numerous practically important settings. Video games frequently unlock new abilities at level transitions (the player-character in the classic game \textit{Dave} acquires a ``shoot'' action in level~3, having only movement actions---up, left, right---in earlier levels). Robots may be fitted with new end-effectors mid-deployment. Skill-learning agents may discover new affordances through interaction. In all these cases, retraining from scratch wastes the knowledge encoded in the agent's state representations, value estimates, and exploration history.

Existing work addresses closely related but distinct settings. \emph{Growing Action Spaces} (GAS)~\cite{farquhar2020growing} requires a predefined nested action hierarchy $\mathcal{A}_0 \subset \mathcal{A}_1 \subset \dots \subset \mathcal{A}_{N-1}$ that must be specified before training begins---it cannot handle truly unexpected additions. \emph{Morphin}~\cite{delarosa2026morphin} detects action-space changes via Page-Hinkley drift and expands Q-tables on the fly, but is tabular only and does not scale to deep reinforcement learning with function approximation. \emph{Headless-AD}~\cite{sinii2023headless} uses action embeddings with an InfoNCE loss to generalise to unseen action-set sizes at test time, but addresses bandit and MDP settings rather than deep RL with complex state spaces, and does not model incremental introduction during training. \emph{Progressive Neural Networks}~\cite{rusu2016progressive} prevent catastrophic forgetting by adding frozen columns for new tasks, but incur linear parameter growth and are designed for task boundaries rather than action-space boundaries within a single task.

To the best of our knowledge, no existing deep RL method handles the scenario where a truly unexpected action appears mid-training, without a predefined hierarchy, a tabular state space, or the requirement to retrain from scratch. In this paper we present a design artifact---a reference architecture and validation harness---that targets this gap. Our contributions are:

\begin{itemize}
  \item \textbf{A reference architecture}, Action-Incremental DQN, with a typed \texttt{ModelConfig} contract, a factorized Q-function $Q(s,a) = f_Q(\varphi(s), e_a)$, a dynamically expandable action-embedding table, and a progressive three-phase unfreezing schedule.
  \item \textbf{A validation harness} covering shape correctness (12 tests), gradient flow (5 tests), RL-specific correctness properties (10 tests), numerical stability including bf16 forward-pass safety (9 tests), and synthetic micro-benchmarks (17 tests)---53 automated tests in total.
  \item \textbf{A 14-condition ablation study} characterising the design space along the dimensions of embedding initialisation, exploration bonus, freeze schedule, Q-function form, auxiliary loss, embedding dimension, and similarity metric.
  \item \textbf{A reference implementation} included in the accompanying artifacts, with a profiler, ablation runner, and demoware (CartPole with simulated action expansion) for reproducible follow-up work.
\end{itemize}

Crucially, we do \emph{not} claim empirical superiority over baselines on real-world tasks. The custom Gymnasium benchmark environment, trained-model evaluation, and baseline comparisons (retrain-from-scratch, fixed-action oracle) are explicitly out of scope for this paper and identified as the natural next step. This paper is appropriate for a workshop or systems-and-reproducibility track: it provides the architectural blueprint, the validation infrastructure, and the honest characterisation of what is---and is not---yet known.

The remainder of the paper is organised as follows. Section~\ref{sec:related} surveys related work. Section~\ref{sec:method} describes the architecture and its design rationale. Section~\ref{sec:setup} details the validation harness. Section~\ref{sec:results} presents synthetic micro-benchmark and ablation results. Section~\ref{sec:discussion} discusses implications and limitations. Section~\ref{sec:conclusion} concludes with concrete follow-up directions.


\section{Related Work}
\label{sec:related}

We organise related work along four axes: predefined action hierarchies, on-the-fly tabular expansion, in-context action generalisation, and task-incremental continual learning.

\subsection{Predefined Action Hierarchies}

Farquhar et al.~\cite{farquhar2020growing} proposed Growing Action Spaces (GAS), which extends DQN to settings where the action set grows according to a known nested hierarchy $\mathcal{A}_0 \subset \mathcal{A}_1 \subset \dots \subset \mathcal{A}_{N-1}$. The core insight is the \emph{monotonicity property}: the optimal value function satisfies $V^*_i(s) \leq V^*_j(s)$ for $i < j$, ensuring that larger action spaces never degrade optimal value. GAS achieves this through parent-child Q-decomposition, where each action's Q-value is expressed relative to its parent in the hierarchy. However, the hierarchy must be fully specified before training begins, and the curriculum of action introductions is hand-designed. This precludes the unexpected-addition scenario that motivates our work. More recently, Seyde et al.~\cite{seyde2024growing} proposed Growing Q-Networks (GQN), which extends the hierarchical principle to continuous action spaces but similarly requires prior knowledge of the action structure.

\subsection{On-the-Fly Tabular Expansion}

De la Rosa et al.~\cite{delarosa2026morphin} introduced morphin, which uses Page-Hinkley drift detection to identify when the action space has changed and expands the Q-table by inserting new rows. The method is self-adaptive, adjusting the exploration rate $\varepsilon$ and learning rate $\alpha$ after each detected change. While morphin demonstrates that on-the-fly detection and expansion is feasible, it operates exclusively in the tabular setting and does not generalise to the deep RL regime where the state space is continuous and function approximation is required. The authors explicitly identify that their drift detector fails when the new reward distribution is a subset of previously seen values.

\subsection{In-Context Action Generalisation}

Sinii et al.~\cite{sinii2023headless} proposed Headless-AD, which decouples the action representation from the policy architecture by using random action embeddings at training time and an InfoNCE loss to align action semantics. At test time, the method can handle more actions than were seen during training by embedding new actions into the same space. This demonstrates that action embeddings can generalise across action-set sizes, which provides evidence for a key assumption of our architecture. However, Headless-AD operates in bandit and MDP settings (not deep RL with complex observations), and crucially does not address the incremental introduction of actions \emph{during training}---it only handles generalisation to larger fixed sets at test time.

\subsection{Task-Incremental Continual Learning}

Progressive Neural Networks~\cite{rusu2016progressive} (PNN)---along with progressive learning for classification~\cite{venkatesan2016progressive}---address catastrophic forgetting in task-incremental learning by allocating a new network column for each new task, with lateral connections to previously learned columns. While PNN prevents forgetting by construction, it incurs linear parameter growth in the number of tasks, requires task identity at test time, and adds entire network columns rather than lightweight action-head expansions. Elastic Weight Consolidation~\cite{kirkpatrick2017overcoming} and other regularisation-based continual learning methods offer an alternative, but are designed for task boundaries rather than action-space boundaries within a single task. The dueling network architecture~\cite{wang2016dueling} provides a complementary decomposition that separates state-value estimation from action advantages, which our dueling variant adopts as a potential mitigation for action-space-induced value shift.

\subsection{Research Gap}

The literature reveals a clear gap: to the best of our knowledge, no existing deep RL method handles the truly unexpected, mid-training introduction of new actions without a predefined hierarchy (GAS), a tabular state space (morphin), or relaxation to a bandit/MDP setting (Headless-AD). The core architectural question---whether a factorized Q-function with learnable action embeddings can provide zero-shot transfer and sub-linear fine-tuning for unexpected actions---remains open. Our work provides a concrete architectural hypothesis and the validation infrastructure to test it.


\section{Method}
\label{sec:method}

\subsection{Architecture Overview}

We present a factorized action-value architecture that decouples action identity from the Q-function. The core design is illustrated in Figure~\ref{fig:arch}.

\begin{figure}[h]
  \centering
  \begin{minipage}{0.85\linewidth}
  {\small
  \begin{verbatim}
         Observation s
                |
                v
    +------------------------+
    | State Encoder phi(s)   |   MLP or CNN
    | Output: R^{256}        |
    | + LayerNorm            |
    +-----------+------------+
                |
                | phi(s)
      +---------+---------+
      |                   |
      v                   v
  +-----------+   +-----------------+
  | Value     |   | Action          |
  | Head V(s) |   | Embedding Table |
  | (dueling) |   | E in R^{|A|x64} |
  +-----+-----+   +--------+--------+
        |                  |
        |           e_a = E[a]
        |                  |
        +--------+---------+
                 |
                 v
        +----------------------+
        | Q-Head f_Q           |
        | MLP(concat[phi; e_a]) |
        | 320 -> 128 -> 1      |
        +----------+-----------+
                   |
                   v
             Q(s, a) = V(s) + A(s,a)  (dueling)
          or Q(s, a) = f_Q(phi, e_a)  (factorized)
  \end{verbatim}
  }
  \end{minipage}
  \caption{Architecture overview. The state encoder $\varphi(s)$ produces a fixed-dimensional representation. Each action $a$ is represented by a learnable embedding $e_a$ in an expandable table. The Q-function $f_Q$ operates on the concatenation $[\varphi(s); e_a]$, decoupling action identity from value estimation.}
  \label{fig:arch}
\end{figure}

\paragraph{State encoder.} The encoder $\varphi: \mathbb{R}^{d_\text{obs}} \to \mathbb{R}^{d_\text{enc}}$ maps observations to a fixed-dimensional representation. For vector observations (default: $d_\text{obs}=64$), $\varphi$ is a 3-layer MLP with hidden dimension 256, ReLU activations, and LayerNorm on the output. For image observations, $\varphi$ uses a standard DQN-style CNN with three convolutional layers followed by a linear projection to $d_\text{enc}=256$.

\paragraph{Action embedding table.} Each action $a$ is associated with a learnable embedding vector $e_a \in \mathbb{R}^{d_\text{emb}}$, stored in a table $E \in \mathbb{R}^{|\mathcal{A}| \times d_\text{emb}}$ (default $d_\text{emb}=64$). The table supports dynamic row insertion via \texttt{add\_embedding()}, which appends a new row without disrupting existing embeddings. Gradient freezing is implemented at the per-row level through a gradient hook that zeroes gradients for specified rows---this is \texttt{torch.compile}-compatible and avoids the need for per-parameter \texttt{requires\_grad} manipulation.

\paragraph{Factorized Q-function.} The Q-function is defined as:

\begin{equation}
Q(s,a) = f_Q([\varphi(s); e_a]), \label{eq:factorized}
\end{equation}

where $f_Q: \mathbb{R}^{d_\text{enc} + d_\text{emb}} \to \mathbb{R}$ is a small MLP (two hidden layers of dimension 128). In the \emph{dueling factorized} variant~\cite{wang2016dueling}, we additionally decompose:

\begin{equation}
Q(s,a) = V(\varphi(s)) + A([\varphi(s); e_a]), \label{eq:dueling}
\end{equation}

where $V$ is a state-value head and $A$ is an advantage head. The advantage values for known actions are mean-centred to mitigate the softmax denominator shift that occurs when new actions are added.

\subsection{Action Expansion Protocol}

When a new action $a_\text{new}$ is detected (the environment's action space grows beyond the agent's current count), the agent follows a four-step protocol:

\paragraph{Step 1: $k$-NN embedding initialisation.} The new action's embedding $e_\text{new}$ is initialised via $k$-nearest-neighbour centroid interpolation in the existing embedding space. Let $E_\text{known} \in \mathbb{R}^{|\mathcal{A}_\text{known}| \times d_\text{emb}}$ be the embeddings of all known actions. We compute:

\begin{align}
\text{centroid} &= \frac{1}{|\mathcal{A}_\text{known}|} \sum_{a \in \mathcal{A}_\text{known}} e_a, \label{eq:centroid} \\
\text{neighbours} &= \text{top-}k_{e \in E_\text{known}} \operatorname{sim}(\text{centroid}, e), \label{eq:knn} \\
e_\text{new} &= \frac{1}{k} \sum_{e \in \text{neighbours}} e. \label{eq:new_emb}
\end{align}

Similarity $\operatorname{sim}$ is cosine similarity by default; Euclidean distance is supported as an alternative. Without any semantic descriptor of the new action, we use the centroid of all known embeddings as a least-committal heuristic initialisation.

\paragraph{Step 2: Optimistic exploration bonus.} The new action receives an additive exploration bonus $\beta$ (default $\beta=2.0$) that decays geometrically with each visit:

\begin{equation}
\beta \leftarrow \max(\beta_\text{min}, \beta \cdot \gamma_\text{bonus}), \quad \gamma_\text{bonus}=0.99.
\end{equation}

This encourages systematic exploration of the new action beyond uniform $\varepsilon$-greedy randomness. A Double DQN variant~\cite{vanhasselt2016double} could be adopted to mitigate potential Q-value overestimation from the additive bonus, but we leave this to future work.

\paragraph{Step 3: Progressive freeze schedule.} To mitigate catastrophic forgetting of existing knowledge, we apply a three-phase unfreezing schedule:

\begin{itemize}
  \item \textbf{Phase 1} (steps $[0, N_1)$, default $N_1=5{,}000$): Only the new action's embedding is trainable. The encoder, Q-head, value head, and all old action embeddings are frozen.
  \item \textbf{Phase 2} (steps $[N_1, N_2)$, default $N_2=10{,}000$): The Q-head and value head are unfrozen. The encoder remains frozen.
  \item \textbf{Phase 3} (steps $[N_2, \infty)$): All parameters are unfrozen with a reduced learning rate ($\text{LR} \leftarrow 0.1 \times \text{LR}_\text{init}$).
\end{itemize}

\paragraph{Step 4: Target network synchronisation.} The target action-embedding table is synchronised from the online network immediately after expansion, ensuring the new embedding is available in the target network. The full target network otherwise follows the standard periodic hard-sync schedule (every $1{,}000$ steps).

\subsection{Training Objective}

The agent is trained using the standard DQN off-policy objective. For a transition $(s, a, r, s', d)$, the TD target is:

\begin{equation}
y = r + \gamma \cdot \max_{a' \in \mathcal{A}_\text{current}} Q'(s', a'), \label{eq:td_target}
\end{equation}

where $\gamma=0.99$ and $Q'$ is the target network. Critically, the max is taken over the \emph{current} action set $\mathcal{A}_\text{current}$, which may be larger than the set that existed when the transition was recorded. This is justified by the monotonicity property of nested action spaces~\cite{farquhar2020growing}: $V^*_i(s) \leq V^*_j(s)$ for $i < j$, so the expanded max correctly reflects the increased value of $s'$.

An optional auxiliary dynamics loss shapes the embedding space:

\begin{equation}
\mathcal{L}_\text{aux} = \text{MSE}\big(\Delta_\text{pred},\ \varphi(s') - \varphi(s)\big), \quad \Delta_\text{pred} = \text{DynamicsHead}([\varphi(s); e_a]),
\end{equation}

weighted by $\lambda_\text{aux}=0.01$ in the total loss $\mathcal{L} = \mathcal{L}_\text{TD} + \lambda_\text{aux} \mathcal{L}_\text{aux}$.

\subsection{Design Rationale}

Table~\ref{tab:design} summarises the key architectural decisions and their justifications.

\begin{table}[h]
  \centering
  \caption{Design decisions and inductive bias justifications.}
  \label{tab:design}
  \begin{tabular}{lp{0.6\linewidth}}
    \toprule
    \textbf{Decision} & \textbf{Justification} \\
    \midrule
    Factorized $Q(s,a) = f_Q(\varphi(s), e_a)$ &
    Decouples action identity from Q-function; intended to support zero-shot Q-values for unseen actions via embedding similarity. \\
    \addlinespace
    Action embedding table &
    Discrete actions lack natural feature vectors; learned embeddings capture similarity structure. \\
    \addlinespace
    $k$-NN centroid initialisation &
    Without semantic descriptors, the centroid of known embeddings is used as a least-committal heuristic initialisation. \\
    \addlinespace
    Progressive freeze (Phase 1$\to$2$\to$3) &
    Isolates new-action parameters during initial learning; state representations transfer perfectly (state space unchanged). \\
    \addlinespace
    Off-policy replay (DQN) &
    Old transitions remain valid after expansion; action indices are stable and the TD target naturally includes new actions. \\
    \addlinespace
    Optional auxiliary dynamics loss &
    Hypothesised to shape embedding space to encode transition dynamics, potentially improving $k$-NN initialisation quality. \\
    \bottomrule
  \end{tabular}
\end{table}

\subsection{Configuration}

All hyperparameters are exposed through a typed \texttt{ModelConfig} dataclass. The default configuration (used for all experiments in this paper) is shown in Table~\ref{tab:config}.

\begin{table}[h]
  \centering
  \caption{Default \texttt{ModelConfig} hyperparameters.}
  \label{tab:config}
  \begin{tabular}{lrc}
    \toprule
    \textbf{Parameter} & \textbf{Default} & \textbf{Description} \\
    \midrule
    $\texttt{obs\_dim}$ / $\texttt{encoder\_dim}$ & 64 / 256 & Observation / encoder dimension \\
    $\texttt{d\_embedding}$ & 64 & Action embedding dimension \\
    $\texttt{k\_nn}$ & 3 & Neighbours for $k$-NN initialisation \\
    $\texttt{q\_form}$ & factorized & Q-function variant \\
    $\texttt{optimistic\_init\_bonus}$ & 2.0 & Exploration bonus for new actions \\
    $\texttt{freeze\_encoder\_steps}$ & 10{,}000 & Phase 1--2 duration (env steps) \\
    $\texttt{freeze\_q\_head\_steps}$ & 5{,}000 & Phase 2 entry (env steps) \\
    $\texttt{expansion\_cooldown}$ & 1{,}000 & Min steps between expansions \\
    $\gamma$ / $\texttt{lr}$ & 0.99 / $3\times10^{-4}$ & Discount factor / learning rate \\
    Total parameters & 272{,}513 & All components (default config) \\
    \bottomrule
  \end{tabular}
\end{table}


\section{Validation Setup}
\label{sec:setup}

Our validation harness is designed to characterise the architecture along four layers, following standard practice for reproducible ML systems research. We emphasise that this harness evaluates \emph{structural and numerical} properties of the architecture, not task performance on real environments. All evaluations use synthetic proxy data (random states, simulated transitions) unless otherwise noted.

\paragraph{Layer 1a: Shape correctness.} Twelve tests verify that every component produces the expected output dimensions for both vector and image observations, for single and batched inputs, and before and after action-space expansion. The batch-computed Q-values are verified against per-action computation (consistency at $10^{-5}$ tolerance).

\paragraph{Layer 1b: Gradient flow.} Five tests verify that gradients flow through all trainable parameters, that no NaN gradients are produced, that the embedding gradient-freeze mask correctly zeroes gradients for frozen rows, and that in Phase~1 only the new embedding receives gradients.

\paragraph{Layer 1c: RL-specific correctness.} Ten tests verify that actions fall within the valid range (before and after expansion), that the exploration entropy is adequate under high $\varepsilon$, that the target network remains stable between synchronisation intervals, that the expansion cooldown guard blocks rapid successive expansions, and that the exploration bonus decays monotonically.

\paragraph{Layer 1d: Numerical stability.} Nine tests verify that no component produces NaN or Inf values in forward passes under standard conditions, extreme input values ($\times 10^4$), and after action expansion. The encoder and Q-head forward passes are verified to produce finite output in bf16 precision, confirming mixed-precision readiness.

\paragraph{Layer 2: Synthetic micro-benchmarks.} Seventeen tests characterise domain-specific properties:
\begin{itemize}
  \item \textbf{Q-degradation metric:} After expansion with Phase~1 freezing, old-action Q-values are preserved to tolerance $10^{-6}$ (no training steps).
  \item \textbf{Monotonicity property:} The maximum Q-value over the expanded action set is at least as large as over the original set, consistent with Farquhar et al.~\cite{farquhar2020growing}.
  \item \textbf{Multiple expansions:} Forward pass, Q-values, and TD updates remain valid through three sequential expansions ($2\to3\to4\to5$ actions).
  \item \textbf{Trainability:} A 200-step CartPole rollout with simulated action expansion produces finite TD losses throughout.
\end{itemize}

\paragraph{Layer 3: Ablation study.} Fourteen single-field \texttt{ModelConfig} changes are evaluated on a proxy metric: Q-value statistics (mean, std, range) computed over 50 random states, plus TD loss on a synthetic batch and the $Q$-value of a newly added action immediately after expansion. The baseline configuration is evaluated alongside each ablation. All ablations share the same random seed.

\paragraph{Layer 4: Profiling.} Forward pass, training step, and action-expansion cost are profiled using \texttt{torch.profiler}. FLOPs are estimated as $2\times$ parameter count for forward pass and $6\times$ for forward+backward.

\paragraph{Compute environment.} All tests and ablations were run on CPU (no GPU available in the validation CI). The total runtime of the full 53-test suite is under 30 seconds. The ablation runner completes in under 60 seconds. Profiling estimates forward pass time at ${<}\,1$ms per step and training step at ${<}\,5$ms per step (batch size 64).

\paragraph{What the harness does not check.} Our harness explicitly does \emph{not} measure: (a) task performance on real environments, (b) cumulative reward curves, (c) comparison against baselines (retrain-from-scratch, fixed-action oracle), or (d) statistical significance across multiple seeds. These are identified as required next steps in Section~\ref{sec:discussion}.

\subsection{Demoware: CartPole with Simulated Expansion}

For development and integration testing, we provide \texttt{ExpandingActionWrapper}, a Gymnasium wrapper that simulates action-space expansion at a configurable step threshold by injecting a new action (mapped to a no-op) into a \texttt{Discrete} action space. This is demoware only---the new action has no semantic effect in CartPole---but it validates that the agent's expansion protocol executes correctly in an end-to-end environment loop.


\section{Results}
\label{sec:results}

\subsection{Validation Harness}

All 53 tests in the validation harness pass. Table~\ref{tab:tests} summarises the coverage.

\begin{table}[h]
  \centering
  \caption{Validation harness test coverage. All 53 tests pass.}
  \label{tab:tests}
  \begin{tabular}{lcc}
    \toprule
    \textbf{Test class} & \textbf{Tests} & \textbf{Coverage area} \\
    \midrule
    \texttt{TestShapes} & 12 & Component shapes, batch consistency, expansion \\
    \texttt{TestGradients} & 5 & Gradient flow, NaN, freeze mask, Phase~1 isolation \\
    \texttt{TestRLProperties} & 10 & Action range, exploration, target sync, bonus, cooldown \\
    \texttt{TestNumerics} & 9 & NaN/Inf, bf16, extreme values, TD loss finiteness \\
    \texttt{TestRLBenchmarks} & 8 & Q-degradation, monotonicity, rollout sanity, expansions \\
    \texttt{TestSyntheticBenchmarks} & 9 & Discounted return, TD target, $k$-NN, params, CNN \\
    \midrule
    Total & 53 & \\
    \bottomrule
  \end{tabular}
\end{table}

\subsection{Synthetic Micro-Benchmarks}

Table~\ref{tab:benchmarks} presents the key synthetic benchmark results. These characterise architectural properties under controlled conditions; they are \emph{not} task-performance numbers.

\begin{table}[h]
  \centering
  \caption{Synthetic micro-benchmark results.}
  \label{tab:benchmarks}
  \begin{tabular}{lp{0.4\linewidth}l}
    \toprule
    \textbf{Property} & \textbf{Observation} & \textbf{Condition} \\
    \midrule
    Q-degradation (no training) & Old-action Q preserved at $10^{-6}$ & Phase~1 frozen params \\
    Q-degradation (20 train steps) & Q-norm change ${<}\,5\times$ & Post-expansion training \\
    Monotonicity & $\max Q(A_\text{expanded}) \geq \max Q(A_\text{orig})$ & Random init, no training \\
    Batch Q consistency & Batch vs.\ per-action Q match at $10^{-5}$ & Full action set \\
    TD loss finiteness & Always finite & Pre- and post-expansion \\
    bf16 stability & No NaN/Inf in forward pass & Encoder and QHead \\
    Freeze schedule & Phase 1$\to$2$\to$3 transitions correct & Verified via \\
    & & \texttt{maybe\_update\_freeze\_schedule} \\
    Multiple expansions & All forward/backward ops valid & $2 \to 3 \to 4 \to 5$ actions \\
    \bottomrule
  \end{tabular}
\end{table}

\subsection{Ablation Study}

Table~\ref{tab:ablations} presents the 14-condition ablation study. All results use proxy metrics (Q-value statistics on random states, TD loss on synthetic batches) and should be interpreted as characterising the design space rather than as comparative validation.

\begin{table}[h]
  \centering
  \caption{Ablation study results. All configurations share the same random seed. The proxy metric (Q-stats on random states) cannot substitute for training-based evaluation.}
  \label{tab:ablations}
  \small
  \begin{tabular}{lccccc}
    \toprule
    \textbf{Ablation} & \textbf{Config $\Delta$} & $\mathbf{\bar{Q}}$ & $\mathbf{\sigma_Q}$ & $\mathbf{\mathcal{L}_\text{TD}}$ & $\mathbf{Q(a_\text{new})}$ \\
    \midrule
    Baseline & --- & 0.116 & 0.0022 & 0.846 & 0.032 \\
    No $k$-NN init & $k_\text{nn}=3 \to 0$ & 0.116 & 0.0022 & 0.825 & NaN \\
    No bonus & bonus $2.0 \to 0.0$ & 0.116 & 0.0022 & 0.863 & 0.122 \\
    No freeze & steps $10\text{K} \to 0$ & 0.116 & 0.0022 & 0.824 & 0.086 \\
    Dueling Q & $\texttt{q\_form}=\text{dueling}$ & $-0.178$ & 0.0022 & 0.694 & $-0.711$ \\
    No aux loss & aux $=\text{False}$ & 0.116 & 0.0022 & 0.798 & $-0.109$ \\
    No emb freeze & freeze $=\text{False}$ & 0.116 & 0.0022 & 0.807 & 0.037 \\
    Small emb & $d_\text{emb}=64\to 8$ & 0.046 & 0.0005 & 0.760 & $-0.271$ \\
    Large emb & $d_\text{emb}=64\to 256$ & $-0.160$ & 0.0027 & 0.902 & $-0.201$ \\
    Euclidean sim & cosine$\to$euclidean & 0.116 & 0.0022 & 0.838 & 0.046 \\
    High bonus & bonus $2.0\to 10.0$ & 0.116 & 0.0022 & 0.780 & 0.072 \\
    Short freeze & steps $10\text{K}\to 100$ & 0.116 & 0.0022 & 0.794 & $-0.053$ \\
    Long freeze & steps $10\text{K}\to 100\text{K}$ & 0.116 & 0.0022 & 0.799 & $-0.046$ \\
    Small Q-head & hidden $128\to 16$ & 0.145 & 0.0014 & 0.831 & 0.002 \\
    Deep Q-head & layers $2\to 4$ & $-0.085$ & 0.0004 & 0.894 & $-0.068$ \\
    \bottomrule
  \end{tabular}
\end{table}

\paragraph{Provisional patterns.} Several structural observations emerge from the ablation data. These observations are \emph{descriptive artifacts of the proxy metric only}---none should be interpreted as evidence of task-performance improvement without training-based validation:

\begin{enumerate}
  \item \textbf{Embedding capacity matters.} Reducing $d_\text{emb}$ from 64 to 8 collapses the Q-value range ($\sigma_Q$ drops from 0.0022 to 0.0005), suggesting that the embedding space lacks the capacity to differentiate actions. This is a potential warning sign for the $k$-NN initialisation hypothesis: if embeddings collapse, the $k$-NN centroid becomes uninformative.
  \item \textbf{Dueling variant shifts the loss landscape.} The dueling decomposition produces a lower mean Q and lower TD loss, likely because the centering mechanism subtracts the known-action advantage mean. Whether this translates to better or worse task performance is unknown.
  \item \textbf{Auxiliary loss changes the sign of $Q(a_\text{new})$ on the proxy metric.} Without the auxiliary dynamics loss, the new action's post-expansion Q-value is negative ($-0.109$), compared to positive for the baseline ($0.032$). This is a purely descriptive observation about the proxy metric: Q-values on random states are not evidence of initialization quality or embedding semantics. Whether the sign change reflects better-initialised embeddings, or is simply an artefact of random-state statistics, can only be resolved through task-based evaluation.
  \item \textbf{Proxy metrics cannot distinguish $k$-NN from random init.} The baseline and \texttt{no\_knn\_init} conditions produce identical Q-stats on random states, as expected. This reinforces that the core hypothesis---$k$-NN initialisation provides better-than-random zero-shot Q-values---can only be tested through training on an actual RL task.
\end{enumerate}


\section{Discussion}
\label{sec:discussion}

\subsection{Interpretation}

Our results validate the \emph{structural} properties of the Action-Incremental DQN architecture. The factorized Q-function correctly computes Q-values before and after expansion; the gradient-freeze mask isolates new-action parameters in Phase~1; the $k$-NN initialisation produces finite, non-NaN embeddings; the TD loss remains finite throughout expansion; and bf16 forward passes are numerically stable. These are necessary conditions for the architecture to be useful, but they are not sufficient to establish empirical effectiveness.

The ablation study reveals design-space structure that would inform future work: embedding dimension is a critical capacity parameter (small $d_\text{emb}$ collapses the action representation); the auxiliary dynamics loss changes the sign of $Q(a_\text{new})$ on the proxy metric---this is a purely descriptive bookkeeping observation and should not be taken as evidence that the loss improves embedding quality or task utility; and the dueling variant fundamentally changes the loss landscape. All of these observations require training-based validation before any conclusions about task performance can be drawn.

\subsection{Comparison to Prior Work}

On structural grounds, our architecture differs from GAS~\cite{farquhar2020growing} in that it does not require a predefined hierarchy---new actions are integrated via $k$-NN in embedding space rather than through a parent-child Q-decomposition. It differs from morphin~\cite{delarosa2026morphin} by operating in the deep RL regime with function approximation. It differs from Headless-AD~\cite{sinii2023headless} by addressing incremental introduction during training rather than generalisation to fixed larger sets at test time. Whether these structural differences translate to practical advantages remains an open empirical question.

\subsection{Limitations}
\label{sec:limitations}

We explicitly enumerate the limitations of this work:

\begin{enumerate}
  \item \textbf{No trained-model evaluation.} All results are from synthetic proxy metrics (random-state Q-statistics, simulated transitions). No RL agent was trained on an actual task with reward feedback.
  \item \textbf{No custom benchmark environment.} The DAVE-game analogue, continuous control, and goal-conditioned scenarios described in the research contract all raise \texttt{NotImplementedError} in our codebase. The CartPole demoware maps new actions to a no-op, which is degenerate.
  \item \textbf{No baseline comparisons.} Retrain-from-scratch, fixed-action-oracle, zero-initialisation, and GAS-adapted baselines are not yet implemented. The sub-linear fine-tuning and no-forgetting claims cannot be validated without them.
  \item \textbf{Single-seed statistics.} All tests and ablations use a single random seed. The research contract requires $\geq\!10$ seeds with mean $\pm$ std reporting.
  \item \textbf{Embedding quality uncharacterised.} No t-SNE visualisation, cosine similarity matrix analysis, or $k$-NN initialisation-quality metric has been computed. The hypothesis that the auxiliary dynamics loss produces ``semantically meaningful'' embeddings remains untested.
  \item \textbf{Softmax denominator shift unmeasured.} The dueling centering mechanism is implemented as a mitigation for denominator shift, but the magnitude of the shift has not been empirically characterised.
  \item \textbf{Single-task focus.} The architecture is evaluated only on a single expansion scenario; cross-task transfer, backwards compatibility after multiple expansions, and sensitivity to expansion order are unstudied.
  \item \textbf{Exploration bonus overestimation.} The additive optimistic bonus may cause Q-value overestimation that propagates through TD backups. Double DQN~\cite{vanhasselt2016double} and bonus clamping are noted as mitigations but not implemented in the current reference implementation.
\end{enumerate}

These limitations are not shortcomings of the implementation---they represent the honest scope of a reference-architecture paper that provides the design, the harness, and the reproducible starting point, leaving empirical validation to subsequent work.


\section{Conclusion}
\label{sec:conclusion}

We presented Action-Incremental DQN, a reference architecture and comprehensive validation harness for the understudied problem of unexpected action-space expansion in deep reinforcement learning. The architecture introduces three mechanisms---factorized Q-function with learnable action embeddings, $k$-NN centroid initialisation for new actions, and progressive three-phase unfreezing---that together form a concrete hypothesis to be investigated: whether deep RL agents can incorporate new actions without catastrophic forgetting or retraining from scratch. Our validation harness provides 53 automated tests covering shape correctness, gradient flow, numerical stability, and synthetic micro-benchmarks, alongside a 14-condition ablation study that characterises the design space.

Three concrete follow-up directions emerge from this work:

\begin{enumerate}
  \item \textbf{Full training-based evaluation.} The immediate priority is building the custom Gymnasium benchmark environment (DAVE-game analogue, continuous control, goal-conditioned affordances) and running training experiments with 10+ seeds per condition, comparing against retrain-from-scratch and fixed-action-oracle baselines.
  \item \textbf{Embedding space analysis.} A systematic characterisation of the learned embedding space---t-SNE visualisation, cosine similarity matrices, and $k$-NN initialisation quality metrics---would test the core assumption that action embeddings encode semantically meaningful structure.
  \item \textbf{Ablation validation on real tasks.} The provisional ablation patterns identified in Section~\ref{sec:results} (embedding capacity, auxiliary loss effect, dueling landscape shift) should be validated on actual RL tasks, potentially with hyperparameter sensitivity sweeps across $k_\text{nn}$, $d_\text{emb}$, and freeze duration.
\end{enumerate}

Our reference implementation and validation harness are included in the accompanying artifact package as a reproducible starting point for the community to build on. The architecture, the harness, and the honest characterisation of what remains unknown are, we hope, a useful contribution towards deep RL systems that can gracefully adapt to changing action repertoires in the wild.


% Acknowledgments (omitted in anonymous submission)


% References
\section*{References}
{\small
\begin{thebibliography}{99}

\bibitem{mnih2015human}
Mnih, V., Kavukcuoglu, K., Silver, D., Rusu, A.A., Veness, J., Bellemare, M.G., Graves, A., Riedmiller, M., Fidjeland, A.K., Ostrovski, G., \& others.
\newblock Human-level control through deep reinforcement learning.
\newblock {\it Nature}, 518(7540):529--533, 2015.

\bibitem{lillicrap2015continuous}
Lillicrap, T.P., Hunt, J.J., Pritzel, A., Heess, N., Erez, T., Tassa, Y., Silver, D., \& Wierstra, D.
\newblock Continuous control with deep reinforcement learning.
\newblock {\it arXiv preprint arXiv:1509.02971}, 2015.

\bibitem{schulman2017proximal}
Schulman, J., Wolski, F., Dhariwal, P., Radford, A., \& Klimov, O.
\newblock Proximal policy optimization algorithms.
\newblock {\it arXiv preprint arXiv:1707.06347}, 2017.

\bibitem{farquhar2020growing}
Farquhar, G., Gustafson, L., Lin, Z., Whiteson, S., Usunier, N., \& Synnaeve, G.
\newblock Growing action spaces.
\newblock In {\it Proceedings of the 37th International Conference on Machine Learning (ICML)}, 2020.

\bibitem{delarosa2026morphin}
de la Rosa, R., Dusparic, I., \& Cardozo, N.
\newblock Adapting the behavior of reinforcement learning agents to changing
  action spaces and reward functions.
\newblock {\it arXiv preprint arXiv:2601.20714}, 2026.

\bibitem{sinii2023headless}
Sinii, V., Nikulin, A., Kolesnikov, S., \& Burda, Y.
\newblock In-context reinforcement learning with variable action spaces.
\newblock In {\it NeurIPS 2023 Workshop on Generalization in Reinforcement Learning}, 2023.

\bibitem{rusu2016progressive}
Rusu, A.A., Rabinowitz, N.C., Desjardins, G., Soyer, H., Kirkpatrick, J., Kavukcuoglu, K., Pascanu, R., \& Hadsell, R.
\newblock Progressive neural networks.
\newblock {\it arXiv preprint arXiv:1606.04671}, 2016.

\bibitem{kirkpatrick2017overcoming}
Kirkpatrick, J., Pascanu, R., Rabinowitz, N., Veness, J., Desjardins, G., Rusu, A.A., Milan, K., Quan, J., Ramalho, T., Grabska-Barwinska, A., \& others.
\newblock Overcoming catastrophic forgetting in neural networks.
\newblock {\it Proceedings of the National Academy of Sciences}, 114(13):3521--3526, 2017.

\bibitem{venkatesan2016progressive}
Venkatesan, R. \& Er, M.J.
\newblock A novel progressive learning technique for multi-class classification.
\newblock {\it Neurocomputing}, 207:310--321, 2016.

\bibitem{seyde2024growing}
Seyde, T., Schwarting, W., Karaman, S., \& Rus, D.
\newblock Growing Q-networks: learning continuous control with growing action spaces.
\newblock In {\it Proceedings of the 6th Conference on Robot Learning (CoRL)}, PMLR 242, 2024.

\bibitem{wang2016dueling}
Wang, Z., Schaul, T., Hessel, M., van Hasselt, H., Lanctot, M., \& de Freitas, N.
\newblock Dueling network architectures for deep reinforcement learning.
\newblock In {\it Proceedings of the 33rd International Conference on Machine Learning (ICML)}, 2016.

\bibitem{vanhasselt2016double}
van Hasselt, H., Guez, A., \& Silver, D.
\newblock Deep reinforcement learning with double Q-learning.
\newblock In {\it Proceedings of the 30th AAAI Conference on Artificial Intelligence}, 2016.

\end{thebibliography}
}


% Appendix
\appendix
\section{Reproducibility Instructions}

Run all commands from the repository root (the parent directory of \texttt{coder/}, \texttt{validator/}, and \texttt{paper/}):

\subsection{Validation Harness}
\begin{verbatim}
# Smoke tests (22 tests, no GPU required)
PYTHONPATH=$PWD/coder:$PYTHONPATH python coder/smoke_test.py

# Comprehensive test suite (53 tests)
PYTHONPATH=$PWD/coder:$PYTHONPATH pytest validator/test_model.py -v --tb=short
\end{verbatim}

\subsection{Ablation Study}
\begin{verbatim}
PYTHONPATH=$PWD/coder:$PYTHONPATH python validator/run_ablations.py
\end{verbatim}

\subsection{Profiling}
\begin{verbatim}
# Forward pass profile
PYTHONPATH=$PWD/coder:$PYTHONPATH python validator/profile_model.py \
    --mode forward --steps 20

# Training step profile
PYTHONPATH=$PWD/coder:$PYTHONPATH python validator/profile_model.py \
    --mode train --steps 10

# Action expansion profile
PYTHONPATH=$PWD/coder:$PYTHONPATH python validator/profile_model.py \
    --mode expansion

# Component parameter counts
PYTHONPATH=$PWD/coder:$PYTHONPATH python validator/profile_model.py \
    --mode component_params
\end{verbatim}

\subsection{Demoware (CartPole with Simulated Expansion)}
\begin{verbatim}
PYTHONPATH=$PWD/coder:$PYTHONPATH python coder/train.py \
    --env cartpole --n_episodes 2000
\end{verbatim}

\subsection{Dependencies}
\begin{itemize}
  \item Python $\geq$ 3.9, PyTorch $\geq$ 2.0, Gymnasium $\geq$ 0.28, NumPy
  \item No GPU required for any test, ablation, or profiling command
\end{itemize}


% NeurIPS Checklist
\newpage
\section*{NeurIPS Paper Checklist}

\paragraph{1. Claims.}
Do the main claims made in the abstract and introduction accurately reflect the paper's contributions and scope?
\textbf{Yes.} The paper claims only to present a reference architecture and validation harness for the problem of unexpected action-space expansion. All contributions are phrased in terms of design artifacts and validated structural properties, not empirical superiority. The absence of trained-model evaluation and real-dataset results is explicitly stated throughout.

\paragraph{2. Experiments.}
Does the paper provide sufficient information to reproduce the experiments?
\textbf{Yes.} The validation harness produces deterministic results with a fixed seed. Full reproducibility instructions, including exact commands for all 53 tests, the 14-condition ablation runner, and profiling scripts, are provided in the appendix. All hyperparameters are documented in the \texttt{ModelConfig} dataclass with defaults listed in Table~\ref{tab:config}.

\paragraph{3. Theory.}
Does the paper make theoretical contributions and, if so, are they supported?
\textbf{Partially.} The paper leverages the monotonicity property of nested action spaces (Farquhar et al., 2020) to justify the TD target formulation. The $k$-NN centroid initialisation is motivated as a least-committal heuristic. However, no new theoretical results (theorems, proofs) are presented.

\paragraph{4. Datasets.}
Does the paper use datasets and, if so, are they documented?
\textbf{No.} The paper does not use any real-world datasets. All evaluations use synthetic random states or the CartPole-v1 environment from Gymnasium (for demoware purposes only).

\paragraph{5. Code.}
Does the paper release code and, if so, is it documented?
\textbf{Yes.} The reference implementation is included in the accompanying artifacts. It includes:
\begin{itemize}
  \item \texttt{model\_config.py}: typed dataclass for all hyperparameters
  \item \texttt{networks.py}: all network components with docstrings and type annotations
  \item \texttt{agent.py}: ActionIncrementalDQN agent with complete public API documentation
  \item \texttt{env\_wrapper.py}: ExpandingActionWrapper for simulated expansions
  \item \texttt{train.py}: training loop with CLI interface
  \item \texttt{smoke\_test.py}: 22-test smoke suite
  \item \texttt{test\_model.py}: 53-test pytest suite
  \item \texttt{run\_ablations.py}: 14-condition ablation runner
  \item \texttt{profile\_model.py}: profiling scripts
\end{itemize}

\paragraph{6. Broader impacts.}
Does the paper discuss potential broader impacts?
\textbf{Minimal.} This paper presents a reference architecture and validation harness for a foundational RL capability (adapting to changing action spaces). The work is at an early, pre-empirical stage. Potential societal impacts are commensurate with those of deep RL generally and are not amplified by this specific contribution. In application domains such as robotics and game AI where agents may acquire new capabilities online, the safety implications of acting on newly added actions before their effects are well-characterised should be carefully considered: a robot fitted with a new end-effector or a game agent that unlocks a novel ability could cause unintended outcomes if the new action is deployed before the value estimate is calibrated. Poorly calibrated exploration of new actions---for example, an additive optimistic bonus that causes sustained Q-value overestimation---could lead to unsafe behaviour in safety-critical settings. Conversely, the ability to incorporate new capabilities without retraining could reduce the computational cost and carbon footprint of lifelong learning systems. These considerations are preliminary and would benefit from domain-specific risk analysis in follow-up work. No immediate dual-use concerns are identified.

\end{document}
